\section{Related Work}
\label{sec:related-work}

\paragraph{Efficient Evaluation.}

Prior work reduces evaluation cost by estimating complete-benchmark results
from item subsets. Fixed-form methods use cross-model correlations, IRT, clustering, or benchmark-level optimization
\citep{vivek2024anchor,maiapolo2024tinybenchmarks,perlitz2024efficient,kipnis2025metabench,fogliato-etal-2024-precise}.
Other methods tailor subsets to the target model, optimize ranking preservation,
or jointly select observed items and predict unobserved outcomes
\citep{yuan2025beyond,saranathan2025sublime,li2025active}. Adaptive evaluation,
building on computerized adaptive testing \citep{lord1980applications}, selects
questions by their expected information about the current system
\citep{truong2025reliable}. These methods are typically calibrated on response
matrices of distinct language models. We instead compare historical caching,
reusable fixed subsets, and adaptive testing on temporally ordered evaluations
of one evolving production agent.



\paragraph{Evaluating LLM Agents}
Agent evaluation is especially expensive because each item may require a full
trajectory involving repository interaction, browsing, tool calls, or
multi-turn interaction
\citep{jimenez2024swe,mialon2024gaia,yao2024tau}. 
Such work develops benchmarks and infrastructure
for comparing agents, while our setting concerns recurring evaluations generated as one agent evolves.


\paragraph{Continuous Evaluation of Evolving Systems}
Our setting also relates to regression-test minimization, selection, and
prioritization for evolving software
\citep{rothermel1997safe,yoo2012regression}. Continuous-integration methods use
code changes and historical executions to allocate testing resources
\citep{kim2002history,elbaum2014techniques,machalica2019predictive}, while
non-deterministic agent outcomes resemble flaky tests
\citep{luo2014empirical}. Conventional regression testing seeks faults using
signals such as coverage or change impact, but changes to an LLM agent do not
directly identify affected benchmark questions. We therefore use historical
question-level outcomes and a partial current run to estimate the complete
benchmark score and run ranking.


